\section{Toward a Competitive Analysis Under Warm-Pool Coupling}
\label{sec:theory}

The empirical result of \S\ref{sec:evaluation:lechowicz} ---
that the provably-optimal-competitive-ratio Lechowicz algorithm
performs \emph{worst} of all baselines on FaaS workloads --- demands
a theoretical explanation. This section makes a partial attempt:
\emph{(i)} we formalize the cost model that distinguishes the FaaS
scheduling problem from the classical one-way trading framework, and
\emph{(ii)} we state the only competitive lower bound we have proven
(the $\beta=0$ specialization, Theorem~\ref{thm:online-lower-bound}),
which simply recovers the classical one-way trading bound and
therefore does not by itself explain the empirical Lechowicz failure
mode.

The honest scope: we do \emph{not} prove a strong lower bound for
Lechowicz on the full warm-pool-coupled cost. The natural
N-invocation adversarial construction does not establish such a bound
in the single-container model, because the idle energy of a single
shared container cannot be charged per-pending-invocation without
double-counting (\S\ref{sec:theory:why_hard} works through this in
detail). We therefore state our finding as a \emph{conjecture},
supported by the empirical evidence of
\S\ref{sec:evaluation:lechowicz} but not yet proven. The challenge of
proving the conjecture, and what its resolution would require, is
itself a useful contribution.

\subsection{The Warm-Pool-Coupled Cost Model}
\label{sec:theory:model}

We restrict to a single region (Lechowicz's setting) so that all
spatial confounders are out of scope and the analysis is about
temporal scheduling alone. The simulator of \S\ref{sec:simulator}
permits a per-(region, function) warm pool of arbitrary size, but
for the formal analysis we present two model variants:
\emph{single-container} (Variant S, matching Lemma~\ref{lemma:tradeoff})
and \emph{independent-pending} (Variant P, an idealization that captures
the per-invocation idle costs that arise during deferral). The
distinction matters: the empirical failure mode of Lechowicz arises
from Variant-P-like effects (pending deferrals each holding resources)
that the simulator implements implicitly through its capacity-bounded
warm pool, but that the strict single-container model does not
capture.

\paragraph{Common inputs.} A sequence of invocations
$\mathcal{I} = \langle i_1, \ldots, i_N \rangle$ arriving online
at times $t_1 \le \cdots \le t_N$. Each invocation has identical
execution time $\tau_e$, cold-start duration $\tau_c$, active power
$P_a$, idle power $P_i$, SLA deadline $D \ge \tau_e$. Carbon intensity
$C: \mathbb{R}_{\ge 0} \to [L, U]$ is revealed causally. Warm-pool
TTL $T_w$ is platform-given.

\paragraph{Decisions.} Each algorithm chooses start times
$s_k \in [t_k, t_k + D - \tau_e]$; we write $c_k = s_k + \tau_e$ for
completion.

\paragraph{Variant S (single-container cost).}
A single warm container is shared across all invocations. The
container is active during $[s_k, c_k]$ (executing $i_k$), idle
during $[c_k, \min(s_{k+1}, c_k + T_w)]$, and expired thereafter
unless reused. Total cost is
\[
\mathrm{Cost}_S = \sum_{k=1}^{N}\!\Big[P_a \tau_e C(s_k)
 + \mathbf{1}[i_k\text{ cold}] P_a \tau_c C(s_k)
 + P_i \cdot \mathrm{IdleTime}_k \cdot \bar{C}_k\Big]
\]
with $\mathrm{IdleTime}_k = \min(s_{k+1} - c_k, T_w)$ and
$\bar{C}_k$ the mean intensity over the idle interval.
By convention, $\mathrm{IdleTime}_0 = 0$ (no pre-$i_1$ container) and
$s_{N+1} = \infty$ so $\mathrm{IdleTime}_N = T_w$.

\paragraph{Variant P (independent-pending cost).}
Each invocation $i_k$ is associated with a dedicated container
that is provisioned (idle) from $t_k$ until $s_k$, then active from
$s_k$ to $c_k$, then idle from $c_k$ until $c_k + T_w$ (TTL) or until
reused. This captures the regime where deferring an invocation
\emph{holds resources} during deferral, which is what production
FaaS platforms with non-trivial warm pools do under load. Cost is
\[
\mathrm{Cost}_P = \sum_{k=1}^{N}\!\Big[
  \underbrace{P_i (s_k - t_k) \bar{C}_k^{\mathrm{pre}}}_{\text{pre-commit idle}}
+ P_a \tau_e C(s_k) + \mathbf{1}[\text{cold}] P_a \tau_c C(s_k)
+ \underbrace{P_i \cdot T_w^{\mathrm{eff}}_k \cdot \bar{C}_k^{\mathrm{post}}}_{\text{post-exec idle}}\Big]
\]
where $T_w^{\mathrm{eff}}_k$ is the post-execution idle period (similar
to Variant~S).

The simulator of \S\ref{sec:simulator} implements something
intermediate: warm containers are pooled per (region, function), and
the per-(region, function) pool is shared across concurrent
invocations of that function in that region. Under heavy load
(arrival rate $\gg 1/T_w$), the pool behaves more like Variant~P;
under light load, more like Variant~S. The reported simulator results
(Tables~\ref{tab:headline}--\ref{tab:realboth_headline}) therefore
do \emph{not} cleanly correspond to either theoretical variant; they
reflect the production-realistic regime where idle-energy is charged
when warm containers outlive their next reuse.

\subsection{The $\beta = 0$ Lower Bound}
\label{sec:theory:beta0}

Under either variant, when $P_i = 0$ (no idle-energy term), the cost
reduces to $\sum_k P_a \tau_e C(s_k)$ plus cold-start terms, which is
the classical one-way trading problem
\citep{el-yaniv2001,lechowicz2024online}.

\begin{theorem}[Classical one-way trading lower bound]
\label{thm:online-lower-bound}
In the restricted cost model with $P_i = 0$, every deterministic
online algorithm $\mathcal{A}$ operating purely on the temporal axis
(no spatial routing) has competitive ratio
$\mathrm{CR}(\mathcal{A}) \ge \sqrt{r}$, where $r = U/L$.
\end{theorem}

\begin{proof}[Proof sketch]
The classical argument from \citet{el-yaniv2001} applies directly.
The adversary plays $C(t) = \sqrt{LU}$ until $\mathcal{A}$ commits;
if $\mathcal{A}$ commits at $\sqrt{LU}$, the adversary drops $C$ to
$L$ (OPT would have paid $L$, ratio $\sqrt{r}$); if $\mathcal{A}$
defers, the adversary raises $C$ to $U$ ($\mathcal{A}$ commits at
$U$, OPT paid $\sqrt{LU}$, ratio $\sqrt{r}$). With $P_i = 0$ the
idle term is irrelevant.
\end{proof}

This is a sanity-check rather than a strong result. It says nothing
about why Lechowicz fails empirically, since Lechowicz \emph{achieves}
$\sqrt{r}$ in the $\beta = 0$ regime.

\subsection{Why the Warm-Pool-Coupled Lower Bound is Hard}
\label{sec:theory:why_hard}

A natural attempt is to prove $\mathrm{CR}(\mathcal{A}_{\text{Lech}})
= \Omega(\beta)$ in Variant~S, using a long-running adversarial trace
that holds $C = U$ for $T_w$ seconds before dropping to $L$. We sketch
why this attempt fails in Variant~S and what changes in Variant~P.

\paragraph{Single-container attempt (Variant~S) — does not work.}
Consider $N$ invocations arriving at $t_k = k\Delta_t$ with
$\Delta_t = T_w/N$ and $C(t) = U$ for $t < T_w$, $C(t) = L$ thereafter.

Under Lechowicz on Variant~S: all $N$ invocations defer to $t = T_w$
when $C$ drops to $L$, then execute back-to-back warm. The pre-$i_1$
container's idle period is at most $T_w$ at $C = U$ (single physical
container, by Variant-S definition). The inter-execution idle periods
($c_k \to s_{k+1}$ for $k=1,\ldots,N-1$) are approximately zero
(back-to-back execution after $T_w$). The post-execution trailing
idle is $T_w$ at $C = L$.
\[
\mathrm{Cost}_S(\mathcal{A}_{\text{Lech}}) \le P_i T_w U + N P_a \tau_e L + P_i T_w L.
\]

Under OPT on Variant~S (commit each at arrival, at $C = U$): inter-arrival
idle periods $\Delta_t$ at $C = U$ between consecutive completions,
then trailing $T_w$ at $C = L$.
\[
\mathrm{Cost}_S(\mathrm{OPT}) \approx N P_a \tau_e U + (N-1) P_i \Delta_t U + P_i T_w L
 = N(P_a \tau_e U) + P_i T_w U (1 - 1/N) + P_i T_w L.
\]

The ratio
$\mathrm{Cost}_S(\mathcal{A}_{\text{Lech}}) / \mathrm{Cost}_S(\mathrm{OPT})$
does \emph{not} grow with $N$ or with $\beta$: both costs have a
$\Theta(P_i T_w \cdot \text{const})$ leading idle term, and the
execution costs differ by a factor of $r = U/L$ amortized over $N$
invocations. The asymptotic ratio (as $N \to \infty$) is
$O(1) + O(1/r)$, trivially close to 1.

\emph{The naive summation $\sum_k P_i (T^\star - t_k) U$, which
attributes idle cost per-pending-invocation, over-counts in
Variant~S: the same physical container is charged $N$ times for one
contiguous idle period.} The construction is therefore not a valid
lower bound in Variant~S.

\paragraph{Independent-pending attempt (Variant~P) — gives the bound.}
In Variant~P, each $i_k$ has its own pending container that idles
from $t_k$ to $s_k$. Under Lechowicz, pending container for $i_k$
idles $(T_w - t_k) = (N-k)\Delta_t$ seconds at $C = U$:
\[
\mathrm{Cost}_P(\mathcal{A}_{\text{Lech}}) = \sum_{k=1}^{N} P_i (T_w - k\Delta_t) U
 + N P_a \tau_e L + N P_i T_w L
 \approx \frac{P_i T_w U N}{2} + N P_a \tau_e L + N P_i T_w L.
\]
Under OPT on Variant~P (commit-at-arrival), each container's
pre-commit idle is zero, post-commit idle $T_w$:
\[
\mathrm{Cost}_P(\mathrm{OPT}) = N P_a \tau_e U + N P_i T_w \bar{C}^{\mathrm{post}}.
\]
The amortized per-invocation ratio is
\[
\frac{\mathrm{Cost}_P(\mathcal{A}_{\text{Lech}})/N}{\mathrm{Cost}_P(\mathrm{OPT})/N}
\to \frac{P_i T_w U / 2 + P_a \tau_e L + P_i T_w L}{P_a \tau_e U + P_i T_w \bar{C}^{\mathrm{post}}}
\]
which approaches $\Theta(\beta)$ if $\bar{C}^{\mathrm{post}}$ is small
relative to $U$ (i.e., if the post-commit window is at low carbon),
and a constant otherwise. The bound is much weaker than the naive
$\beta + 1/r$ a per-pending-invocation accounting would suggest, and
it depends on the adversary's freedom to choose the post-commit
carbon trace.

\paragraph{The simulator's regime.} The simulator implements neither
Variant~S nor Variant~P exactly: warm containers are pooled per
(region, function), so a pending $i_k$ does not provision a new
container if a compatible warm one exists. Under heavy load, the pool
size grows toward the number of concurrent pending invocations
(approaching Variant~P); under light load, a single warm container
serves successive arrivals (approaching Variant~S). Lechowicz's
empirical $-336\%$ vs FIFO on real data
(Table~\ref{tab:lechowicz}) is consistent with the intermediate
regime: not the trivial Variant~S worst case, not the
$\Theta(\beta)$ Variant~P worst case.

\subsection{Conjecture and Open Problems}
\label{sec:theory:conjecture}

We state what we believe but cannot yet prove.

\begin{conjecture}
\label{conj:lech-lb}
There exists a multi-container (Variant~P-like) cost model that
captures FaaS production behavior, in which any algorithm of the
Lechowicz form (threshold-based deferral without an idle-energy
charge) has competitive ratio $\Omega(\beta)$ in the regime
$\beta \gg \sqrt{r}$.
\end{conjecture}

Resolving this conjecture requires three pieces of work we have not
done:
\emph{(i)} a formal multi-container cost model that matches FaaS
production semantics (warm-pool sharing, capacity, function-level
isolation),
\emph{(ii)} an adversary construction in that model that withstands
the per-invocation idle accounting (not the double-counted Variant~S
attempt above), and
\emph{(iii)} a matching upper-bound algorithm to show the lower bound
is tight, ideally a randomized variant of \GreenFaaS{}'s break-even
gate that achieves $O(\sqrt{r})$ on stationary inputs.

We view this as the principal theoretical follow-up the paper
suggests. The empirical evidence
(\S\ref{sec:evaluation:lechowicz},
Tables~\ref{tab:lechowicz}, \ref{tab:realboth_headline}) is
consistent with the conjecture but does not prove it.

\subsection{What This Section Establishes}
\label{sec:theory:summary}

The theoretical contribution of this paper is honestly modest:
\emph{(i)} a formalization of the warm-pool-coupled cost model
(\S\ref{sec:theory:model}) that distinguishes single-container,
independent-pending, and intermediate variants, useful for
follow-up work;
\emph{(ii)} the classical one-way trading lower bound
(Theorem~\ref{thm:online-lower-bound}) under the $\beta=0$
specialization, which is a sanity check rather than a strong result;
\emph{(iii)} a careful working-out of why the natural N-invocation
adversarial construction fails in the single-container model
(\S\ref{sec:theory:why_hard}), which clarifies what a future proof
would need to handle;
\emph{(iv)} an open conjecture (\S\ref{sec:theory:conjecture}) that
the warm-pool-coupled regime admits an $\Omega(\beta)$ lower bound on
Lechowicz-style algorithms.

The empirical evidence for the conjecture is strong (Lechowicz
empirically performs worst of all baselines on FaaS workloads,
including those where its theoretical $\sqrt{r}$ guarantee should
apply). A complete theoretical resolution is open work.
